\section{Introduction}

\IEEEPARstart{D}{ecision} trees and Na\"ive Bayes are two of the oldest
workhorses of supervised classification, and each covers the other's
weak spot: a tree splits on whichever attribute looks most informative at
each node but ignores the rest of the joint distribution, while NB uses
every attribute but assumes they are independent given the class. Farid
et al.~\cite{farid2014hybrid} proposed two hybrids that combine the two
models in opposite directions. Algorithm 1 trains a NB classifier, uses it
to flag the training instances it misclassifies as noise, deletes them,
and grows a C4.5 tree on what remains. Algorithm 2 grows a tree first,
turns the depth at which each attribute is tested into a weight
($1/\sqrt{d}$, zero for attributes the tree never tests), and plugs those
weights into NB's class-conditional likelihood. Both hybrids were reported
to outperform their un-hybridized counterpart on ten benchmark
datasets (nine from the UCI repository plus NSL-KDD), and the paper has since accumulated a substantial citation
record built on those numbers.

Three properties of the original evaluation raise concern independently
of whether the numbers are correct. First, the paper never states whether
the noise filter and the attribute weights were fit once on the full
dataset or refit inside each cross-validation fold; if it is the former,
every deleted instance and every attribute weight was chosen with
knowledge of the labels of the instances later used for testing, which is
exactly the kind of preprocessing leakage that inflates reported accuracy
in unrelated
domains~\cite{cawley2010overfitting,demircioglu2021measuring,kapoor2023leakage,bouke2023empirical}.
Second, Algorithm 1 equates ``the instance NB gets wrong'' with ``the
instance is noise,'' a design that later work on label-noise filtering
identifies as one of the weaker choices available, because a single model
is both the source and the arbiter of its own
disagreement~\cite{brodley1999identifying,jiang2024correction,northcutt2021confident}.
Third, the two hybrids were evaluated separately: nobody, including the
original paper, reports what happens when the filter and the weighting
scheme are chained, in either order, inside one pipeline.

This paper reports a replication built to check the first two concerns
directly and to measure the third. Our contributions are:
\begin{itemize}
\item \textbf{A faithful replication (R1)} of Farid et al.'s two hybrids
using the paper's own components (Weka J48 (pruned, default options) as
C4.5 and a textbook NB on the original nominal attribute columns) under
two evaluation protocols: protocol A (fit-before-CV), where the filter or
weights are fit once on the whole dataset, and protocol B
(refit-per-fold), where every stage is refit inside each training fold.
\item \textbf{A quantified leakage audit (E3a)} that isolates the effect of
protocol alone, holding the algorithms, data encoding and classifiers
fixed, and repeats the audit on an independent codebase (sklearn entropy
trees, one-hot attributes) to check the finding is not an artifact of one
implementation.
\item \textbf{A two-order pipeline study (E9)} that chains Algorithm 1 and
Algorithm 2 in both possible orders under protocol B, with a final NB or
decision tree, to measure whether order or combination changes the
honest picture.
\item \textbf{Honest statistical reporting} of every comparison as a
Wilcoxon signed-rank test over the ten datasets, Holm-adjusted, with an
explicit acknowledgment of the low power that ten paired
datasets gives~\cite{demsar2006statistical}.
\end{itemize}
All numbers in this paper are computed by a script from CSV files produced
by runnable code; none are transcribed from earlier notes or hand-copied
from the original paper's tables.
